\subsection{Dynamic Power Estimation}
\label{subsec:dynamic_power}

A fundamental objective of this research is to quantify the energy efficiency gains achieved by transitioning from a continuous \gls{cnn} to a bio-plausible \gls{snn}.
Traditional \gls{cnn} architectures are characterised by static power consumption; they perform a dense, fixed number of \gls{mac} operations for every forward pass, regardless of the information density within the input scene. 
In contrast, \glspl{snn} leverage event-driven dynamics, where metabolic energy is consumed only when a neuron reaches its firing threshold and generates a discrete spike.

The Dynamic Power experiment investigates the reactive scalability of the architecture by mapping instantaneous energy consumption against the temporal density of the input stream. 
This analysis serves to validate the event-driven nature of the \gls{snn}, demonstrating that its metabolic cost is a dynamic function of activation sparsity—mimicking biological neural efficiency by consuming energy only in response to informative stimuli.

\subsubsection{\gls{cnn} Power Estimation}

In a standard \gls{cnn}, the total energy per frame ($E_{CNN}$) is calculated by multiplying the total number of \glspl{flop} by the cost of a 32-bit \gls{mac} ($E_{MAC}$):

\begin{equation}
  E_{CNN} = \sum_{l=1}^{L} (FLOPs_l \times E_{MAC})
\end{equation}

\noindent Where $E_{MAC} = 4.6\text{ pJ}$ per 32-bit floating point operation, as referenced in Table \ref{tab:energy_constants} \cite{horowitz20141}.

The number of \glspl{flop} is the total number of \glspl{mac} operations required to produce the output feature map; in a \gls{cnn}, this number is static, regardless of what the image contains.
The \glspl{flop} of a \gls{cnn} can be calculated as follows:

\begin{equation}
     FLOPs_l = (H_{out} \times W_{out}) \times (K \times K) \times C_{in} \times C_{out}
\end{equation}

Where $(H_{out} \times W_{out})$ is the number of pixels in the output frame, $(K \times K)$ is the size of the convolutional kernel, $C_{in}$ is the number of input channels, and $C_{out}$ is the number of output channels.


\subsubsection{\gls{snn} Power Estimation}

For the spiking architecture, the calculation is scaled by the activation sparsity ($\zeta$) of each layer.
The metabolic cost of the spiking backbone can be quantified via the \gls{sop} metric;
unlike the static complexity of the \gls{cnn} baseline, the \gls{sop} count for a given layer is a dynamic function, which can be calculated as follows:
\begin{equation}
SOPs_l = FLOPs_l \times \zeta \times T
\label{eq:sop_zeta}
\end{equation}

\noindent where $\zeta_l$ represents the average firing rate of layer $l$, $T$ is the number of temporal simulation steps ($T=10$).

In order to obtain $\zeta_l$, a specialised \texttt{SpikeRateMonitor} tool was developed.
The tool attaches a hook to every \gls{lif} inside the \gls{snn} backbone, which listens to the layer during the forward pass and intercepts the binary spike tensor before it moves to the next layer.
For every layer $l$, the \gls{snn} produces a spike tensor $S_l$ with the shape $[T, B, C, H, W]$, where $T = 10$ is the number of timesteps, $B$ is the batch size, and $C, H, W$ are the spatial dimensions.
Since a spike is binary ($1$ for fire, $0$ for rest), $\zeta$ can be calculated by taking the global mean of the entire tensor:

\begin{equation}
    \zeta_l = \frac{1}{N} \sum_{i=1}^{N} S_{l, i}
\end{equation}

To calculate Dynamic Power, $\zeta$ was calculated by running inference over all validation samples and taking the mean firing rate per layer across all samples in the dataset.

% \todo[inline]{Mention SpikeRateMonitor in Evaluation Suite Section and Diagram}

Because \glspl{snn} utilise binary activations, the expensive multiplication part of the \gls{mac} operation is omitted, leaving only a \gls{ac} of the weight to the membrane potential.
The total energy per frame ($E_{SNN}$) is defined as:

\begin{equation}
\label{eq:e_snn}
  E_{SNN} = \sum_{l=1}^{L} (SOPs_l \times E_{AC})
\end{equation}


\noindent where $E_{AC} = 0.9\text{ pJ}$ is the energy cost of a single floating-point addition, as referenced in Table \ref{tab:energy_constants} \cite{horowitz20141}.

